\section{Introduction}
\label{sec:introduction}

Activation steering has emerged as a lightweight method for controlling language
model behavior without fine-tuning~\citep{turner2023actadd,panickssery2024caa,zou2023repe}.
By adding a learned direction vector to a model's hidden states during the forward
pass, practitioners can increase truthfulness, reduce toxicity, or shift stylistic
register. The approach is attractive because it requires no gradient updates,
operates at inference time, and needs only a small set of contrastive examples
to extract a steering direction.

A natural question arises: must steering vectors be applied at every token
position, or can brief application produce lasting effects? Standard methods
such as Contrastive Activation Addition~\citep{panickssery2024caa} inject
vectors continuously, implicitly assuming that the effect would not persist
without reinforcement. Yet no prior work has directly measured how quickly
steering influence dissipates after the vector is removed. Understanding this
decay is important for three reasons. First, practitioners want to minimize
computational overhead by applying vectors only when needed. Second, the decay
rate reveals how the model processes and absorbs perturbations to its internal
representations. Third, the distinction between representational decay (hidden
states reverting to their natural trajectory) and behavioral persistence
(generated tokens maintaining influence through the context) has direct
implications for steering method design.

We provide the first empirical characterization of post-removal steering vector
decay. Using \gptxl (1.5B parameters) and TransformerLens~\citep{nanda2022transformerlens},
we extract steering vectors for three behavioral dimensions---truthfulness,
positive sentiment, and formal style---via contrastive activation addition.
We then apply each vector for the first $K$ tokens of generation, remove it,
and measure hidden state dynamics token-by-token for 20 subsequent tokens.
We compare six conditions: no steering, continuous steering, autoregressive
generation after removal (\AR), teacher-forced generation with the same steered
tokens (\TF), teacher-forced generation with unsteered tokens (\TFclean), and a
random-vector control.

Our central finding is that representational decay is \emph{instantaneous}.
The cosine similarity delta between steered and unsteered hidden states drops
from ${\sim}0.019$ to ${\sim}0.000$ within a single generation step after
steering removal. This result is consistent across all three behaviors, all
layers tested (8, 16, 24, 32), all multiplier strengths (1--5$\times$), and
all steering durations ($K{=}1$ to $K{=}10$). Autoregressive and teacher-forced
conditions produce identical hidden states (Cohen's $d = 0.000$), confirming
that the model retains no hidden state memory of the perturbation. Yet
behavioral effects persist: the KL divergence between the steered model's
output distribution and the unsteered baseline remains elevated for many tokens,
with a half-life exceeding 75 tokens (\figref{fig:behavioral}). This persistence vanishes entirely under
the \TFclean condition, where unsteered tokens replace the steered context,
confirming that behavioral influence is carried by the generated token sequence
rather than by hidden state dynamics.

We make the following contributions:
\begin{itemize}[leftmargin=*,itemsep=0pt,topsep=0pt]
    \item We present the first token-by-token measurement of steering vector
    decay dynamics, showing that representational influence is absorbed within
    a single forward pass.
    \item We decompose steering persistence into representational and behavioral
    components, demonstrating that behavioral effects operate entirely through
    the autoregressive token feedback loop.
    \item We conduct comprehensive ablations across layers, multiplier strengths,
    steering durations, and behavioral dimensions, establishing the universality
    of instantaneous representational decay.
\end{itemize}
